% !TITLE 哨遍·高祖还乡（节选）
% !AUTHOR 睢景臣
% !DYNASTY 元
% !GENRE 散曲
% !ORDER 5

\begin{poem}
社长\textsuperscript{1}排门告示\textsuperscript{2}，但有的差使无推故\textsuperscript{3}，这差使不寻俗\textsuperscript{4}。
一壁厢纳草也根\textsuperscript{5}，一边又要差夫\textsuperscript{6}，索应付\textsuperscript{7}。
又言是车驾\textsuperscript{8}，都说是銮舆\textsuperscript{9}，今日还乡故\textsuperscript{10}。
王乡老\textsuperscript{11}执定瓦台盘\textsuperscript{12}，赵忙郎\textsuperscript{13}抱着酒胡芦\textsuperscript{14}。
新刷来的头巾，恰糨来的绸衫\textsuperscript{15}，畅好是妆幺大户\textsuperscript{16}。
\end{poem}

\begin{annotations}
\notetext{1}{社长：元代乡村基层组织的头目。元代以五十家为一社，设社长。}
\notetext{2}{排门告示：挨家挨户通知。“排门”即挨门挨户。}
\notetext{3}{但有的差使无推故：凡是派到的差使，都不许推脱。但，凡是。推故，找借口推辞。}
\notetext{4}{不寻俗：不寻常、不一般。俗语，指这次差事非同小可。}
\notetext{5}{一壁厢纳草也根：一边要缴纳喂马的草料，还得除根（把草的根部清理干净）。一壁厢，一边。}
\notetext{6}{差夫：差役、徭役。指派老百姓去服劳役。}
\notetext{7}{索应付：必须应付、照办。索，必须。}
\notetext{8}{车驾：皇帝乘坐的车子，代指皇帝。}
\notetext{9}{銮舆：皇帝的车驾。“銮”是车上的铃铛，“舆”是车厢。与“车驾”同义，重复使用显出口语化。}
\notetext{10}{还乡故：回到故乡。“乡故”即故乡，倒装。}
\notetext{11}{王乡老：姓王的乡村长老。“乡老”是乡间有头脸的人物。}
\notetext{12}{瓦台盘：陶制的托盘，用来盛放酒食接待贵客。贫寒人家用瓦制而非瓷器，可见乡村的寒酸。}
\notetext{13}{赵忙郎：姓赵的伙计。“忙郎”是无代口语，指帮工、跑腿的人。}
\notetext{14}{酒胡芦：盛酒的葫芦。村民们拿着简陋的器具去迎接“皇帝”，场面滑稽。}
\notetext{15}{恰糨来的绸衫：刚刚浆洗过的绸衫。糨（jiàng），用米汤或粉浆浸泡衣物使其挺括。乡民为迎接皇帝，刻意打扮。}
\notetext{16}{畅好是妆幺大户：真是装模作样的大户人家。畅好是，真是；妆幺，装模作样。讽刺乡民穿得人模人样，却掩饰不住底层的寒酸。}
\end{annotations}

\begin{appreciation}
《哨遍·高祖还乡》是元代文学里最痛快的一支曲子之一——它把皇帝拉回了人间。哨遍是曲牌名，原是一套"高祖还乡"套曲的头一支，这里节选开头。

刘邦当上皇帝后衣锦还乡，史书里是风光场面，睢景臣偏不写风光，写皇帝还没到，村子先乱了。节选就是这场忙乱：社长挨家挨户传话，说有大人物要来，差使一律不许推脱。一边要交喂马的草料，还得一根一根去掉根；一边又派差役，让你去服徭役。全村人鸡飞狗跳，就为等一队车马。

细节里处处是老百姓的眼睛：姓王的乡老捧着瓦台盘，姓赵的忙郎抱着酒葫芦——瓦台盘是陶的，乡下人家里只有这个；新刷的头巾，刚浆洗的绸衫，穿得整整齐齐，怎么看怎么别扭。"畅好是妆幺大户"——真是装模作样的大户人家。你看，老百姓什么都知道，只是不敢说。这一句"妆幺大户"，刺既扎在乡民身上，也扎在皇帝身上。

节选之外才是高潮：天子仪仗到了，旌旗招展，乡民越看越眼熟——旗帜上画的不就是白环套着迎霜兔吗？再往下，有人认出来了：你不是刘三吗？你当过亭长，喝过酒，赊过账，借过我的粮，还欠我的豆子钱。好好的，改了姓、更了名，唤做汉高祖，就以为没人认得你了！

这才是元曲敢干的事。诗词里写皇帝，要么歌功颂德，要么小心翼翼；元曲不一样，元代文人没了科举这条路，没了敬畏，说话直来直去。当时扬州的曲家都写过这个题目，只有睢景臣的传了下来，因为只有他敢让乡民掀皇帝的底。皇帝的威仪是全村人用两天脚程搭起来的布景，戏台底下站着的还是那个刘三。

哥哥想让你记住的是：不管一个人后来多了不起，村里人都记得他小时候什么样。你将来也是一样——走得多远，回来还是老家这个小孩，邻居还叫你的小名。所以犯不着装，像曲子里的绸衫，浆洗得再挺括，一眼就看穿了。
\end{appreciation}
